\section{Real-Data Application: Rotterdam Breast Cancer}
\label{sec:real}

We evaluate the flow on the Rotterdam breast cancer dataset
\citep{foekens2000rotterdam} ($n = 2{,}982$ patients, $d = 7$ covariates:
age, menopausal status, tumor size, histological grade, number of positive
lymph nodes, progesterone receptor level, and estrogen receptor level).
The endpoint is time to disease recurrence ($T$), with death without
recurrence as the censoring event ($C$). This is a natural setting for
joint $(T, C)$ modeling: death competes with recurrence, so censoring is
plausibly informative.

\paragraph{Protocol.}
We use $5$-fold cross-validation repeated $5$ times ($25$ train/test
splits). All covariates are standardized to zero mean and unit variance.
Times are measured in years, clipped to $[0.01, 15]$. The flow uses
$K = 4$ tanh components, hidden dimension $64$, and the same
hyperparameters as the simulation study. We compare against DSM ($K{=}4$
mixture of LogNormals), multivariate Cox PH (partial likelihood with
Breslow baseline), and Kaplan--Meier. Metrics are the IPCW-weighted
integrated Brier score \citep{graf1999assessment} and Harrell's
concordance index.

\paragraph{Results.}
Table~\ref{tab:rotterdam} reports mean $\pm$ standard deviation across
the $25$ splits.

\begin{table}[t]
\caption{Rotterdam breast cancer ($n = 2{,}982$, $d = 7$): $5 \times 5$
cross-validated IBS and C-index. Best in \textbf{bold}.}
\label{tab:rotterdam}
\centering
\small
\begin{tabular}{lcc}
\toprule
Method & IBS ($\downarrow$) & C-index ($\uparrow$) \\
\midrule
Cond.\ flow (ours) & $\mathbf{0.180\,{\pm}\,0.004}$ & $\mathbf{0.679\,{\pm}\,0.016}$ \\
DSM ($K{=}4$)       & $\mathbf{0.180\,{\pm}\,0.005}$ & $0.679\,{\pm}\,0.018$ \\
Cox PH              & $0.184\,{\pm}\,0.004$ & $0.674\,{\pm}\,0.015$ \\
Kaplan--Meier       & $0.211\,{\pm}\,0.003$ & $0.500\,{\pm}\,0.000$ \\
\bottomrule
\end{tabular}
\end{table}

The flow and DSM are statistically tied (IBS $= 0.180$, C-index
$\approx 0.679$), both outperforming Cox PH on calibration (IBS $0.180$
vs.\ $0.184$) and discrimination (C-index $0.679$ vs.\ $0.674$). All
covariate-aware methods strongly dominate KM ($0.211$ / $0.500$).

Two observations are worth noting. First, the flow achieves competitive
performance on real data with $d = 7$ covariates using the same
architecture and hyperparameters as the $d = 1$ simulations, suggesting
the method scales without scenario-specific tuning. Second, the flow
jointly models event and censoring times, which may provide more
informative predictions in settings where the censoring mechanism itself
is of clinical interest (e.g., death competing with recurrence).
